% GENERATED_BY: oc_core_monograph_translation_draft_pipeline_v1.py
% AI_ASSISTED_TRANSLATION_DRAFT: TRUE
% THEOREM_REVIEW_STATUS: PENDING
% THEOREM_REVIEW_MODE: FORMAL_AI_GATE__CODEX_CLI_CHATGPT
% THEOREM_REVIEW_REVIEWER_ID:
% THEOREM_REVIEW_AUTHORITY_CLASS:
% THEOREM_REVIEWED_AT:
% THEOREM_REVIEW_DOSSIER_REF:
% SOURCE_LANGUAGE: EN
% TARGET_LANGUAGE: RU
% SOURCE_EN_REF: content/k_levels/k3.tex
% ================================================================
% ==== FILE: content/k_levels/k3.tex
% ================================================================

% ==============================
%  Ontology of Continua — Core
%  K-level Module: K3
%  File: content/k_levels/k3.tex
%  Status: FULLY DEFINED
% ==============================

\subsubsection{\texorpdfstring{$K_3$}{K_3} Обзор}
\label{sec:k3-overview}

$K_3$ — это уровень, на котором
\emph{химическая организация} становится
структурно возможной. Он представляет переход от чисто геометрических
и топологических континуумов ($K_0$--$K_2$) к системам, в которых:
\begin{itemize}
    \item реакции, преобразования и конверсии появляются как допустимые
          потоки,
    \item потенциалы соответствуют реакционным affinities и градиентам
          концентраций,
    \item границы $\partial\Omega$ включают молекулярно-композиционные и
          топологические ограничения,
    \item минимальное автокаталитическое замыкание становится представимым,
    \item возникают первые структурно согласованные сети (сети RAF).
\end{itemize}

Размерностный переход $\Psi_{2\to3}$ происходит, когда несовместимые
химические различия не могут быть выражены в рамках осей $\{A_1,A_2\}$
уровня $K_2$. Вводится новая ось $A_3$, представляющая отдельную
размерность
\emph{пространства состава}, которая делает возможным химическое
возникновение континуумов.

$K_3$ служит структурной подосновой для всех последующих биологических и
когнитивных континуумов, но сам по себе не является живым.

% ================================================================
\subsubsection{\texorpdfstring{$\Omega(K_3)$}{\Omega(K_3)} Пространство состояний}

Пространство состояний $K_3$:
\[
\Omega(K_3) = \{ (\mathbf{c}, \mathbf{r}, \Gamma) \mid
\mathbf{c} \in \mathbb{R}_{\ge0}^N,\;
\mathbf{r} \in \mathbb{R}_{\ge0}^M,\;
\Gamma \subseteq R \},
\]
где:
\begin{itemize}
    \item $\mathbf{c}$: вектор концентраций молекулярных видов,
    \item $\mathbf{r}$: вектор реакционных потоков,
    \item $\Gamma$: множество активированных реакций (стехиометрическая структура).
\end{itemize}

Свойства:
\begin{enumerate}
    \item $\Omega(K_3)$ допускает нелинейные преобразования, индуцированные
          реакционными сетями.
    \item Допустимые состояния должны удовлетворять ограничениям баланса масс и
          порогам из $M_3$.
    \item Подсети RAF как структурно согласованные области встраиваются в
          $\Omega(K_3)$.
\end{enumerate}

Стабильность $K_3$ зависит от того, содержит ли $\Omega(K_3)$ какие-либо
замкнутые автокаталитические подмножества с устойчивым потоком.

% ================================================================
\subsubsection{\texorpdfstring{$\partial\Omega(K_3)$}{\partial\Omega(K_3)} Граница}

$\partial\Omega(K_3)$ включает:
\begin{itemize}
    \item состояния, где какая-либо концентрация становится отрицательной
          (запрещено),
    \item расходящиеся потоки, превышающие $\Theta_{\mathrm{flux}}$,
    \item химические конфигурации, нарушающие стехиометрическую реализуемость,
    \item режимы коллапса: pH-коллапс, осмотический коллапс,
          нестабильность реакционных цепочек,
    \item потерю автокаталитического замыкания.
\end{itemize}

Ключевым для $K_3$ является:
\textbf{потеря замыкания} приводит к разрывности $\Omega(K_3)$,
что предотвращает устойчивую организацию.

% ================================================================
\subsubsection{\texorpdfstring{$A(K_3)$}{A(K_3)} Оси}

$K_3$ поддерживает три независимые структурные оси:
\[
A(K_3) = \{A_1, A_2, A_3\},
\]
где $A_3$ является новой осью состава.

$A_3$ кодирует:
\begin{itemize}
    \item стехиометрическую структуру,
    \item степени свободы направления реакций,
    \item композиционные различия, которые нельзя отобразить на
          геометрические или фазовые различия согласованности $K_2$.
\end{itemize}

Независимость:
\[
A_3 \notin \mathrm{span}\{A_1,A_2\}.
\]

Так выражается непреодолимая новизна химической организации.

% ================================================================
\subsubsection{\texorpdfstring{$P(K_3)$}{P(K_3)} Потенциалы}

Химические потенциалы включают:
\[
P(K_3) = \{ \mu_i,\; \Delta G_r,\; P_{\mathrm{affinity}},\;
           P_{\mathrm{mix}},\; P_{\mathrm{exchange}}\},
\]
где:
\begin{itemize}
    \item $\mu_i$: потенциалы по видам,
    \item $\Delta G_r$: термодинамически свободные энергии реакций,
    \item $P_{\mathrm{affinity}}$: движущие силы реакций,
    \item $P_{\mathrm{mix}}$: энтропийные потенциалы смешивания,
    \item $P_{\mathrm{exchange}}$: обмен со средой (
          \emph{если} разрешено $M_3$).
\end{itemize}

Химическая организация возникает, когда эти потенциалы создают
стабильную, самоподдерживающуюся структуру.

% ================================================================
\subsubsection{\texorpdfstring{$\Theta(K_3)$}{\Theta(K_3)} Пороги}

Соответствующие пороги:
\[
\Theta(K_3) =
\{\Theta_{\mathrm{closure}},\;
  \Theta_{\mathrm{flux}},\;
  \Theta_{\mathrm{pH}},\;
  \Theta_{\mathrm{osm}},\;
  \Theta_{\mathrm{grad}},\;
  \Theta_{\mathrm{react}}\}.
\]

Интерпретация:
\begin{itemize}
    \item $\Theta_{\mathrm{closure}}$: минимальные условия для
          автокаталитического замыкания,
    \item $\Theta_{\mathrm{flux}}$: максимально устойчивая скорость реакций,
    \item $\Theta_{\mathrm{pH}}$: допустимая кислотно-основная
          неуравновешенность,
    \item $\Theta_{\mathrm{osm}}$: порог осмотической устойчивости,
    \item $\Theta_{\mathrm{grad}}$: границы градиента концентраций,
    \item $\Theta_{\mathrm{react}}$: энергетический порог реакций.
\end{itemize}

Переход через эти пороги приводит к коллапсу (Секц.~the K3 collapse discussion).

% ================================================================
\subsubsection{\texorpdfstring{$J(K_3)$}{J(K_3)} Потоки}

Допустимые потоки:
\[
J(K_3) = \{\text{реакционные потоки},\; \text{диффузионный транспорт},\;
         \text{потоки обмена},\; \text{автокаталитическое усиление}\}.
\]

Типы качественные:
\begin{itemize}
    \item линейные массовые потоки,
    \item нелинейные автокаталитические потоки,
    \item выравнивающие потоки, предотвращающие самовозрастание,
    \item обменные потоки через протобеременные границы.
\end{itemize}

Автокаталитические потоки центральны: они определяют структурное
усиление, необходимое для жизнеспособности $K_3$.

% ================================================================
\subsubsection{\texorpdfstring{$C(K_3)$}{C(K_3)} Циклы}

Циклы в $K_3$ включают:
\[
C(K_3) = \{\text{реакционные циклы},\;
           \text{автокаталитические петли},\;
           C_{\mathrm{RAF}}\}.
\]

$C_{\mathrm{RAF}}$ — минимальный замкнутый RAF-цикл, удовлетворяющий:
\[
r_i \\in \\text{RAF} \\iff \\text{все реагенты $r_i$ и катализатор\r\nпроизведены внутри данного множества.}
\]

Свойства:
\begin{itemize}
    \item RAF-циклы обеспечивают структурную персистентность,
    \item их функционал мощности $\Pi(C_{\mathrm{RAF}})$ определяет,
          может ли быть организация устойчивой,
    \item существование $C_{\mathrm{RAF}}$ необходимо (но не достаточно)
          для перехода к $K_4$.
\end{itemize}

Когда $C_{\mathrm{RAF}}$ становится устойчивым, $K_3$ поддерживает
прототипы временной упорядоченности.

% ================================================================
\subsubsection{\texorpdfstring{$\tau(K_3)$}{\tau(K_3)} Время}

Время в $K_3$ возникает, когда автокаталитические циклы поддерживают
ненулевую мощность цикла:
\[
\tau(K_3) \text{ существует тогда и только тогда, когда } \Pi(C_{\mathrm{RAF}}) >
\Theta_{\mathrm{time}}.
\]

Временная структура слабее, чем в $K_2$:
\begin{itemize}
    \item время определяется через итерацию реакционных циклов,
    \item время исчезает, если циклы гасят или переходят к нулевому потоку.
\end{itemize}

Стабильная временная упорядоченность — необходимое условие для перехода
к биологическому $K_4$.

% ================================================================
\subsubsection{\texorpdfstring{$k(K_3)$}{k(K_3)} Континуумность}

Применяется общая формула:
\[
k(K_3) =
\frac{\mu(\Omega(K_3))}{\mu(S(K_3))}
\cdot
\frac{|A(K_3)|}{|A(M_3)|}
\cdot
\frac{\sum \mathrm{Stab}(C)}{\sum \mathrm{MaxStab}(C)}
\cdot
(1 - \frac{\sigma(J)}{\sigma_{\max}})
\cdot
(1 - \frac{T}{\Theta_{\mathrm{stab}}})_+.
\]

Интерпретация:
\begin{itemize}
    \item существование RAF-цикла повышает $k(K_3)$,
    \item сильный дисбаланс потоков снижает $k(K_3)$,
    \item осмотический или pH-стресс драматически уменьшает $k(K_3)$,
    \item $k(K_3)=0$ соответствует химической смерти.
\end{itemize}

% ================================================================
\subsubsection{\texorpdfstring{$T(K_3)$}{T(K_3)} Структурное напряжение}

Источники напряжения:
\begin{itemize}
    \item несовместимые реакционные потоки,
    \item сильные концентрационные градиенты,
    \item pH-дисбаланс,
    \item осмотические градиенты,
    \item отсутствующие катализаторы (неполное замыкание).
\end{itemize}

Когда:
\[
T(K_3) > \Theta_{\mathrm{stab}},
\]
автокаталитическая организация коллапсирует.

% ================================================================
\subsubsection{\texorpdfstring{$E(K_3)$}{E(K_3)} Энергия}

Энергетический функционал:
\[
E(K_3) =
\sum_i \mu_i c_i +
\sum_r \Delta G_r\, r +
E_{\mathrm{grad}} + E_{\mathrm{mix}} + E_{\mathrm{osm}}.
\]

В терминах:
\begin{itemize}
    \item химические потенциалы и свободные энергии,
    \item затраты градиентной энергии,
    \item осмотический дисбаланс,
    \item энтропия смешивания.
\end{itemize}

Этот функционал задает химическую устойчивость.

% ================================================================
\subsubsection{\texorpdfstring{$K_3$}{K_3} Операторы (\texorpdfstring{$\Psi$}{\Psi}, \texorpdfstring{$\Phi$}{\Phi}, \texorpdfstring{$\Lambda$}{\Lambda}, \texorpdfstring{$U$}{U}, \texorpdfstring{$\Chi$}{\Chi})}

\begin{itemize}
    \item $\Psi_{3\to4}$ — биологический переход: появление
          компартментации (мембранного закрытия), порождающей $K_4$.
    \item $\Phi$ — эволюция ограничений в $M_3$ (пределы состава,
          допустимость реакций).
    \item $\Lambda$ — оператор стехиометрической согласованности.
    \item $U$ — унификация реакционных подсетей.
    \item $\Chi$ — допускает ветвление в альтернативные химические
          организации (разные RAF-ядра).
\end{itemize}

% ================================================================
\subsubsection{Процессы на \texorpdfstring{$K_3$}{K_3}}

Представительные процессы:
\begin{itemize}
    \item автокаталитическое усиление,
    \item образование перекрестно-каталитических цепей,
    \item pH-дрифт и восстановление,
    \item осмотический дисбаланс и стабилизация,
    \item формирование или распад RAF-подсетей,
    \item структурный дрейф к $K_4$, когда возникает граница $\partial\Omega$.
\end{itemize}

% ================================================================
\subsubsection{Прогнозы для \texorpdfstring{$K_3$}{K_3}}

\begin{enumerate}
    \item Возникновение RAF — пороговый феномен, управляемый
          $\Theta_{\mathrm{closure}}$.
    \item Минимальная химическая организация требует автокаталитических циклов.
    \item pH-коллапс вызывает быстрое разрушение $\Omega(K_3)$.
    \item Кривая осмотического стресса имеет универсальную сигмоидную форму
          ($\Theta_{\mathrm{osm}}$).
    \item Переход к $K_4$ требует формирования полу-устойчивой
          протограницы.
\end{enumerate}

% ================================================================
\subsubsection{Эксперименты для \texorpdfstring{$K_3$}{K_3}}

Эмпирические аналоги:
\begin{itemize}
    \item экспериментальные сети RAF (Steel–Hordijk),
    \item системы полимеризации без протоклеток,
    \item осцилляторы pH-градиентов,
    \item осмотические циклы набухания и разрыва в везикулах без мембран,
    \item автокаталитические реакторы открытых систем.
\end{itemize}

Каждая из них даёт проверки для $\Theta_{\mathrm{closure}}$,
$\Theta_{\mathrm{react}}$ и критериев коллапса.

% ================================================================
\subsubsection{Смерть и коллапс \texorpdfstring{$K_3$}{K_3}}
\label{sec:k3-collapse}

Коллапс происходит, когда:
\[
\Omega(K_3) = \varnothing,
\]
по причине:
\begin{itemize}
    \item разрушения RAF-замыкания,
    \item pH-коллапса,
    \item осмотического разрыва,
    \item неконтролируемого расхода потоков,
    \item несовместимых стехиометрических ограничений.
\end{itemize}

После смерти переход к биологическому $K_4$ невозможен.

% ================================================================
\subsubsection{Фальсифицируемость \texorpdfstring{$K_3$}{K_3}}

Прогнозы, поддающиеся проверке:
\begin{enumerate}
    \item существование минимального порога RAF,
    \item необходимость автокаталитических циклов для устойчивой организации,
    \item воспроизводимая осмотическая оболочка коллапса в разных химических
          средах,
    \item pH-границы, определяющие структурное выживание,
    \item невозможность поддержания химической организации без замыкания.
\end{enumerate}

Нарушения этих предсказаний опровергнут структурную модель.

% ================================================================
\subsubsection{Ветвление / Онтологическая позиция \texorpdfstring{$K_3$}{K_3}}

Ветвление происходит через:
\begin{itemize}
    \item альтернативные RAF-ядра,
    \item различные стехиометрические подсети,
    \item различающиеся режимы обмена с окружающей средой.
\end{itemize}

Позиционно:
\[
K_2 \to K_3 \to K_4,
\]
где $K_3$ является уникальным химическим промежуточным звеном между
физическими и биологическими континуумами.

% ================================================================
\subsubsection{Связь с M-пространствами}

$K_3$ существует внутри $M_3$, который задает:
\begin{itemize}
    \item допустимые диапазоны состава,
    \item допустимые классы реакций,
    \item пределы обмена со средой,
    \item осмотические и pH-ограничения,
    \item структурные условия для возникновения границ $\partial\Omega$.
\end{itemize}

Переход $K_3 \to K_4$ требует от $M_3$ порождения нового допустимого
топологического различения «внутри/снаружи».
